\documentclass[11pt,a4paper]{article}

% ---------------------------------------------------------------- packages
\usepackage[T1]{fontenc}
\usepackage[utf8]{inputenc}
\usepackage{lmodern}
\usepackage[a4paper,margin=2.4cm]{geometry}
\usepackage[dvipsnames,table]{xcolor}
\usepackage{booktabs}
\usepackage{array}
\usepackage{amssymb}
\usepackage{enumitem}
\usepackage{fancyhdr}
\usepackage{parskip}
\usepackage[hidelinks]{hyperref}

\setlist{nosep,leftmargin=1.4em}
\renewcommand{\arraystretch}{1.18}

% ---------------------------------------------------------------- status macros
\newcommand{\yes}{\textcolor{ForestGreen}{$\checkmark$}}
\newcommand{\pp}{\textcolor{orange}{\textbf{$\sim$}}}
\newcommand{\no}{\textcolor{BrickRed}{$\times$}}
\newcommand{\na}{\textcolor{gray}{--}}

% ---------------------------------------------------------------- header / footer
\pagestyle{fancy}
\fancyhf{}
\lhead{\footnotesize Programme-Website Content Review --- Findings Report}
\rhead{\footnotesize\thepage}
\cfoot{\footnotesize Cyprus University of Technology \quad\textbullet\quad For internal use}
\renewcommand{\headrulewidth}{0.4pt}
\renewcommand{\footrulewidth}{0.4pt}

\hypersetup{
  pdftitle={Programme-Website Content Review --- Findings Report},
  pdfauthor={Mihalis Golias, Ioanna Evangelou, Elena Kalotychou},
}

% ================================================================ document
\begin{document}

% ---------------------------------------------------------------- title block
\begin{center}
  {\LARGE\bfseries Programme-Website Content Review}\\[4pt]
  {\large Findings Report}\\[10pt]
  {\normalsize Faculty of Management and Economics\\
  School of Tourism Management, Hospitality and Entrepreneurship\\
  Cyprus University of Technology}\\[10pt]
  {\small Prepared by: Mihalis Golias, Ioanna Evangelou and Elena Kalotychou}\\[2pt]
  {\small Review period: June 2026}
\end{center}

\vspace{4pt}
\hrule
\vspace{10pt}

% ---------------------------------------------------------------- summary box
\noindent\textbf{Purpose.} This report documents the current content of the public programme webpages of the two Faculties, assessed against a consistent set of eight content items. It records \emph{what was found} on the live pages; it does not propose page designs. Its aim is to give each Department a clear, comparable picture of where programme information is complete, partial, or missing, so that gaps can be closed ahead of accreditation review and for prospective-student recruitment.

\tableofcontents
\vspace{6pt}
\hrule

% ================================================================ 1
\section{Executive summary}

Fifteen programme pages were reviewed across four Departments. Overall, the undergraduate and the flagship taught-master pages are well developed, while several postgraduate and all doctoral pages fall short of the standard set by the reference page (the MSc in Shipping). The most important findings are:

\begin{enumerate}[label=\arabic*.]
  \item \textbf{Programme-level Intended Learning Outcomes (ILOs) are missing or only implicit on every page.} No page presents ILOs as an explicit, measurable list; at best they are embedded inside individual course descriptions. This is the single most consistent gap and an explicit accreditation expectation.
  \item \textbf{Two taught-master curricula are not published online at all.} The public pages for the \emph{MSc in Management and AI} and the \emph{MSc in International Tourism and Hospitality Management} function only as gateways: they expose the programme title (and, in one case, a coordinator name) but no aims, course list, credits, descriptions, admission detail, fees, or careers information. These are the highest-priority items.
  \item \textbf{Careers and recognition information is uneven.} Only the Shipping-side pages publish graduate-employment percentages. Most other pages describe sectors and roles but give no statistics. Several programmes are newly established and therefore have no graduate data yet --- a fact that should simply be stated on the page.
  \item \textbf{Doctoral pages are consistently thin.} They generally lack explicit ILOs, a structured curriculum with credits, and graduate-career information; some present only research areas or a current vacancy notice.
  \item \textbf{Course descriptions and credits are incomplete in places.} Several pages list course titles without descriptions, and the two Data Science degrees publish no ECTS values at all.
  \item \textbf{A number of data-integrity inconsistencies were found} (campus, programme title, credit totals, and a duplicated course code), indicating that a content-governance pass is warranted.
\end{enumerate}

% ================================================================ 2
\section{Scope and method}

The review covered the public, prospective-student programme pages of:

\begin{itemize}
  \item \textbf{Faculty of Management and Economics (FME)} --- the Department of Shipping and the Department of Finance, Accounting and Management Science;
  \item \textbf{School of Tourism Management, Hospitality and Entrepreneurship (STMHE)} --- the Department of Management, Entrepreneurship and Digital Enterprise, and the Department of Tourism and Hospitality Management.
\end{itemize}

Each page was assessed against the reference page (the MSc in Shipping) using eight content items grouped under two headings: \emph{academic content} and \emph{outcomes and careers}. Pages were read from the live site during June 2026.

\textbf{Caveat.} Some content on the live site is held behind JavaScript tabs or accordions and may not be retrievable by automated reading; for the two master pages flagged as ``not published online'' this was confirmed through repeated checks of every associated page. Items marked as missing should be spot-checked on the live page before any action is taken.

% ================================================================ 3
\section{What was assessed}

\textbf{Academic content}
\begin{itemize}
  \item \textbf{Aims} --- programme aims and objectives.
  \item \textbf{ILOs} --- explicit, programme-level intended learning outcomes.
  \item \textbf{Course descriptions} --- a description of the content of each course/module.
  \item \textbf{Structure} --- the distribution of courses per semester (the CYQAA ``Table~2'' course distribution).
  \item \textbf{Syllabus} --- depth of course content / syllabus detail.
\end{itemize}

\textbf{Outcomes and careers}
\begin{itemize}
  \item \textbf{Employment} --- employment opportunities and sectors for graduates.
  \item \textbf{Professional recognition} --- professional accreditation or professional-body exemptions, where applicable.
  \item \textbf{Career outcomes} --- graduate destinations: typical job roles and employment statistics.
\end{itemize}

% ================================================================ 4
\section{Summary gap matrix}

\begin{center}
\footnotesize
\setlength{\tabcolsep}{5pt}
\begin{tabular}{@{}p{5.7cm}cccccccc@{}}
\toprule
\textbf{Programme} & \textbf{Aims} & \textbf{ILOs} & \textbf{Desc.} & \textbf{Str.} & \textbf{Syll.} & \textbf{Empl.} & \textbf{Rec.} & \textbf{Car.}\\
\midrule
\multicolumn{9}{@{}l}{\cellcolor{gray!15}\textbf{FME --- Department of Shipping}}\\
MSc in Shipping \textit{(reference standard)} & \yes & \pp & \yes & \yes & \yes & \yes & \yes & \yes\\
BSc in Shipping & \yes & \yes & \yes & \yes & \pp & \yes & \yes & \yes\\
Doctoral Studies in Shipping (PhD) & \yes & \pp & \yes & \yes & \pp & \yes & \yes & \pp\\
\addlinespace[2pt]
\multicolumn{9}{@{}l}{\cellcolor{gray!15}\textbf{FME --- Department of Finance, Accounting and Management Science}}\\
BSc in Commerce, Finance and Shipping \textit{(legacy)} & \yes & \pp & \yes & \yes & \yes & \pp & \yes & \no\\
BSc in Data Science & \yes & \pp & \yes & \yes & \pp & \yes & \no & \yes\\
BSc in Data Science for Economics and Management & \yes & \pp & \yes & \yes & \pp & \yes & \no & \pp\\
BSc in Finance and Accounting & \yes & \pp & \yes & \yes & \yes & \yes & \yes & \pp\\
MSc in Financial Management and Investment & \yes & \no & \pp & \yes & \no & \yes & \yes & \pp\\
PhD in Finance, Accounting and Management Science & \yes & \no & \pp & \yes & \no & \no & \no & \no\\
\addlinespace[2pt]
\multicolumn{9}{@{}l}{\cellcolor{gray!15}\textbf{STMHE --- Dept of Management, Entrepreneurship and Digital Enterprise}}\\
BA in Management, Entrepreneurship and Digital Business & \yes & \pp & \yes & \yes & \yes & \yes & \yes & \pp\\
MSc in Management and AI: Entrep.\ and Digital Transf. & \no & \no & \no & \no & \no & \no & \no & \no\\
PhD in Business Administration and Entrepreneurship & \pp & \no & \no & \no & \no & \no & \na & \no\\
\addlinespace[2pt]
\multicolumn{9}{@{}l}{\cellcolor{gray!15}\textbf{STMHE --- Department of Tourism and Hospitality Management}}\\
Bachelor of Tourism and Hospitality Management & \yes & \pp & \yes & \yes & \yes & \yes & \yes & \pp\\
MSc in International Tourism and Hospitality Management & \no & \no & \no & \no & \no & \no & \no & \no\\
PhD in Hospitality and Tourism Management & \pp & \no & \no & \no & \no & \no & \na & \no\\
\bottomrule
\end{tabular}
\end{center}

\noindent{\footnotesize \textbf{Legend:} \yes~present \quad \pp~partial / implicit \quad \no~missing or not published \quad \na~not applicable.\\
Abbreviations: Desc.~=~course descriptions; Str.~=~structure (semester distribution); Syll.~=~syllabus depth; Empl.~=~employment / sectors; Rec.~=~professional recognition; Car.~=~career outcomes (roles \& statistics).}

% ================================================================ 5
\section{Cross-cutting findings}

\begin{enumerate}[label=\textbf{F\arabic*.},leftmargin=2.6em]
  \item \textbf{Explicit ILOs are absent everywhere.} A single standard ILO block (knowledge / skills / competences) should be adopted and added to every programme page.
  \item \textbf{Two master curricula are not online.} The MSc in Management and AI and the MSc in International Tourism and Hospitality Management publish no curriculum on the live site; their pages act only as gateways.
  \item \textbf{Recognition-and-careers depth is inconsistent.} Only the Shipping-side pages quote employment percentages. Elsewhere, recognition is sometimes only linked, not stated, and career statistics are usually absent. For newly established programmes, the page should state plainly that no graduate data exist yet.
  \item \textbf{Course descriptions versus titles only.} Several pages (the MSc in Financial Management and Investment and both Finance/Management-side doctoral pages in particular) list course titles without descriptions.
  \item \textbf{Missing credits (ECTS).} The two Data Science degrees publish no ECTS values; doctoral taught-course credits are given only as component totals, not per course.
  \item \textbf{Thin doctoral pages.} Doctoral pages generally lack explicit ILOs, a structured curriculum with credits, and graduate-career information.
  \item \textbf{Tuition and fees are rarely shown} on the programme pages themselves.
  \item \textbf{English versions are inconsistent.} Most pages are Greek-first; equivalent English pages are needed for international recruitment and for accreditation review.
  \item \textbf{Stale and duplicate pages.} Discontinued programmes and duplicate, divergent pages remain reachable and should be removed or redirected.
\end{enumerate}

% ================================================================ 6
\section{Findings by programme}

\subsection{FME --- Department of Shipping}

\textbf{MSc in Shipping (reference standard).}\\
\emph{Present:} the benchmark page --- programme aims, a full course list by semester with ECTS, detailed per-course descriptions, professional recognition (Institute of Chartered Shipbrokers exemptions, up to three), and a stated graduate-employment figure.\\
\emph{Gaps:} ILOs are not presented as an explicit, measurable list.

\textbf{BSc in Shipping.}\\
\emph{Present:} explicit aims and learning outcomes, a full eight-semester structure with codes and ECTS, ICS exemptions (up to three), and Department employment data (90.3\% labour-market absorption).\\
\emph{Gaps:} roughly two dozen courses are listed by title only; tuition is not published; ILOs are not a discrete list. The current four-year structure dates from 2024, so the programme has not yet produced its own graduates.

\textbf{Doctoral Studies in Shipping (PhD).}\\
\emph{Present:} aims, a four-stage 240-ECTS structure, admission requirements, and a research-skills focus.\\
\emph{Gaps:} no explicit ILOs; limited per-course detail; graduate/career outcomes are qualitative only, with no statistics.

\subsection{FME --- Department of Finance, Accounting and Management Science}

\textbf{BSc in Commerce, Finance and Shipping (legacy; final intake 2024).}\\
\emph{Present:} a full curriculum with codes and ECTS and professional recognition (ICAEW ACA module exemptions; ACCA).\\
\emph{Gaps:} the majority of courses are titles-only; no employment statistics; ILOs not explicit. \emph{Data integrity:} per-semester ECTS sum to 248 against a stated 240. As the programme is being phased out, the Faculty should decide the extent of any update.

\textbf{BSc in Data Science (first intake 2025).}\\
\emph{Present:} aims, a full course list with translated descriptions for many courses, and the access-framework score; typical graduate roles are listed.\\
\emph{Gaps:} no ECTS values anywhere (per-course or total); no tuition; no named coordinator; professional recognition not addressed; ILOs not explicit. As a new programme it has no graduate statistics yet.

\textbf{BSc in Data Science for Economics and Management (first intake 2025).}\\
\emph{Present:} aims, course structure, and descriptions; typical roles and sectors.\\
\emph{Gaps:} no ECTS; ILOs not explicit; professional recognition not addressed; no graduate statistics (new programme). \emph{Data integrity:} the English title appears inconsistently (``\ldots and Management'' versus ``\ldots and Business'').

\textbf{BSc in Finance and Accounting (first intake 2025).}\\
\emph{Present:} a full curriculum across Finance and Accounting tracks with codes, ECTS, and detailed descriptions; professional recognition (ICAEW, up to twelve exemptions; ACCA, nine).\\
\emph{Gaps:} ILOs not explicit; no employment statistics (new programme). \emph{Data integrity:} some per-semester listings exceed 30 ECTS once the internship/dissertation is included.

\textbf{MSc in Financial Management and Investment.}\\
\emph{Present:} strong careers content (sector and role mapping, on the departmental microsite), recognition framing, a published tuition figure (EUR~5{,}125), and a 90-ECTS structure.\\
\emph{Gaps:} course descriptions are titles-only (no syllabi); ILOs absent. \emph{Data integrity:} one course code (FAMS~566) is used for two different courses.

\textbf{PhD in Finance, Accounting and Management Science.}\\
\emph{Present:} aims, a full taught-course and research-stage list with component ECTS (240 = 48 + 32 + 32 + 128), detailed admission, and research areas with named supervisors.\\
\emph{Gaps:} no per-course ECTS for taught courses; course descriptions titles-only; no ILOs; no employment or career outcomes.

\subsection{STMHE --- Department of Management, Entrepreneurship and Digital Enterprise}

\textbf{BA in Management, Entrepreneurship and Digital Business.}\\
\emph{Present:} a full curriculum with codes and ECTS, descriptions, two specialisations (Entrepreneurship and Innovation; Digital Business), professional recognition (ICAEW/ACCA), and a structured internship.\\
\emph{Gaps:} ILOs not explicit; no employment statistics. \emph{Data integrity:} the page address (Limassol) differs from the campus of record (Paphos); the specialisation name varies (``Digital Business'' versus ``Digital Enterprise'').

\textbf{MSc in Management and AI: Entrepreneurship and Digital Transformation.}\\
\emph{Critical.} The public page is effectively a gateway --- it exposes the programme title, staff, and Departmental research areas only. No published aims, ILOs, course list, ECTS, descriptions, admission detail, tuition, or career information could be found online. Highest-priority remediation.

\textbf{PhD in Business Administration and Entrepreneurship.}\\
\emph{Present:} a general aim, three research areas (Management; Entrepreneurship and Innovation; Information Systems), and staff.\\
\emph{Gaps:} no ILOs; no published curriculum or ECTS; no course descriptions; no career outcomes. Established 2022, with no graduates yet.

\subsection{STMHE --- Department of Tourism and Hospitality Management}

\textbf{Bachelor of Tourism and Hospitality Management.}\\
\emph{Present:} a full curriculum (eight semesters, codes and ECTS), detailed module descriptions, two specialisations (Hospitality; Tourism), professional recognition (ACCA/ICAEW), and a substantial internship.\\
\emph{Gaps:} ILOs not explicit; no employment statistics. \emph{Data integrity:} per-semester ECTS headers do not always match their row totals; two parallel course-code systems are in use; the page address (Limassol) differs from the campus of record (Paphos).

\textbf{MSc in International Tourism and Hospitality Management.}\\
\emph{Critical.} The public page is a gateway exposing only the programme title and a coordinator name. No aims, ILOs, course list, ECTS, descriptions, admission, tuition, or career content could be found online. Highest-priority remediation.

\textbf{PhD in Hospitality and Tourism Management.}\\
\emph{Present:} the live page functions as a doctoral-vacancy notice --- it carries detailed admission requirements and two advertised research themes with supervisors.\\
\emph{Gaps:} no ILOs; no programme structure or ECTS; no course descriptions; no stable research-areas list; no career outcomes.

% ================================================================ 7
\section{Data-integrity issues observed}

These are inconsistencies noted while reading the pages; each should be confirmed and corrected by the responsible Department.

\begin{itemize}
  \item \textbf{Campus.} The BA in Management and the Bachelor of Tourism and Hospitality Management pages give a Limassol address, whereas the campus of record for both is Paphos.
  \item \textbf{Programme title.} The BSc in Data Science for Economics is rendered as both ``\ldots and Management'' and ``\ldots and Business''.
  \item \textbf{Credit totals.} The BSc in Commerce, Finance and Shipping per-semester ECTS sum to 248 against a stated 240; some Finance-and-Accounting and Tourism semesters list more than 30 ECTS once internship/dissertation credits are included.
  \item \textbf{Duplicate course code.} On the MSc in Financial Management and Investment, the code FAMS~566 is used for two different courses.
  \item \textbf{Parallel code systems.} The Bachelor of Tourism and Hospitality Management uses two different course-code schemes for the same courses across its structure and module pages.
\end{itemize}

% ================================================================ 8
\section{Recommendations}

\textbf{Priority~1 --- close the largest gaps}
\begin{itemize}
  \item Publish the full curricula for the MSc in Management and AI and the MSc in International Tourism and Hospitality Management on the live site.
  \item Add an explicit, programme-level ILO block (knowledge / skills / competences) to every programme page.
\end{itemize}

\textbf{Priority~2 --- bring pages up to the reference standard}
\begin{itemize}
  \item Standardise a ``Recognition and Careers'' block on every page: (a) professional recognition/exemptions, or an explicit statement that none apply; (b) typical job roles and sectors; (c) employment statistics, or a clear note that the programme is newly established and has no graduate data yet.
  \item Add course descriptions wherever only titles are shown, and publish ECTS for the Data Science degrees and for doctoral taught courses.
  \item Strengthen the doctoral pages with ILOs, a structured curriculum with credits, and graduate-career information.
\end{itemize}

\textbf{Priority~3 --- consistency and governance}
\begin{itemize}
  \item Provide an equivalent English version of every programme page.
  \item Remove or redirect pages for discontinued programmes and any duplicate, divergent pages.
  \item Correct the data-integrity issues listed in Section~7, and assign a single content owner and a common page template per programme to prevent recurrence.
\end{itemize}

% ================================================================ appendix
\section*{Appendix --- pages reviewed}
\addcontentsline{toc}{section}{Appendix --- pages reviewed}
\footnotesize

\textbf{FME --- Department of Shipping}
\begin{itemize}
  \item MSc in Shipping (reference): \url{https://www.cut.ac.cy/faculties/fme/shipping/degrees/postgraduate/msc_shipping/}
  \item BSc in Shipping: \url{https://www.cut.ac.cy/faculties/fme/shipping/degrees/undergraduate-programmes/}
  \item Doctoral Studies in Shipping: \url{https://www.cut.ac.cy/faculties/fme/shipping/degrees/doctoral.studies/}
\end{itemize}

\textbf{FME --- Department of Finance, Accounting and Management Science}
\begin{itemize}
  \item BSc in Commerce, Finance and Shipping: \url{https://www.cut.ac.cy/faculties/fme/cfs/degrees/undergraduate-programmes/emporiou_xrim_nautilias/}
  \item BSc in Data Science: \url{https://www.cut.ac.cy/faculties/fme/cfs/degrees/undergraduate-programmes/data_science/}
  \item BSc in Data Science for Economics and Management: \url{https://www.cut.ac.cy/faculties/fme/cfs/degrees/undergraduate-programmes/data_science_econ_manag/}
  \item BSc in Finance and Accounting: \url{https://www.cut.ac.cy/faculties/fme/cfs/degrees/undergraduate-programmes/finance/}
  \item MSc in Financial Management and Investment: \url{https://www.cut.ac.cy/faculties/fme/cfs/degrees/postgraduate/msc-finance-business-admin/} (and \url{https://www.cutfams.com/})
  \item PhD in Finance, Accounting and Management Science: \url{https://www.cut.ac.cy/faculties/fme/cfs/degrees/doctoral.studies/}
\end{itemize}

\textbf{STMHE --- Department of Management, Entrepreneurship and Digital Enterprise}
\begin{itemize}
  \item BA in Management, Entrepreneurship and Digital Business: \url{https://www.cut.ac.cy/faculties/tohe/program-in-management/degrees/undergraduate_programmes/}
  \item MSc in Management and AI: \url{https://www.cut.ac.cy/faculties/tohe/program-in-management/degrees/postgraduate/}
  \item PhD in Business Administration and Entrepreneurship: \url{https://www.cut.ac.cy/faculties/tohe/program-in-management/degrees/doctoral_studies/}
\end{itemize}

\textbf{STMHE --- Department of Tourism and Hospitality Management}
\begin{itemize}
  \item Bachelor of Tourism and Hospitality Management: \url{https://www.cut.ac.cy/faculties/tohe/thm/study-guide/undergraduate-programmes/}
  \item MSc in International Tourism and Hospitality Management: \url{https://www.cut.ac.cy/faculties/tohe/thm/study-guide/postgraduate/}
  \item PhD in Hospitality and Tourism Management: \url{https://www.cut.ac.cy/faculties/tohe/thm/study-guide/phdHTM/}
\end{itemize}

\end{document}
